% ============================================================================
\section{Discusión}

Los hallazgos obtenidos permiten responder con claridad a las cuatro interrogantes que guiaron la investigación y contrastarlos con la literatura sobre datos macroeconómicos en tiempo real.

En primer lugar, respecto del carácter sistemático de las revisiones, los resultados confirman lo planteado por \textcite{croushore2011} y el análisis entre nueva información y ruido de medición \parencite{mankiw1986}: a horizonte continuo de once meses (\(n=350\)), las primeras publicaciones del PBI peruano no actúan como pronósticos racionales insesgados. Muestran una subestimación promedio positiva y robusta de \(+0{,}347\) puntos porcentuales (\(p=5{,}3\times 10^{-8}\)), persistencia anual detectable en tiempo real (\(\hat{\rho}_{12}=+0{,}386\)), colas pesadas y rechazo de eficiencia de \textcite{mincer1969} debido al intercepto positivo constante, patrón que persiste incluso después de retirar los cruces de cambio de año base de 2000 y 2014.

En segundo lugar, en cuanto al desempeño de los referentes clásicos, los modelos muestran un comportamiento claramente distinto antes y después de la pandemia. En el tramo previo a 2020 (\(n=190\)), la regresión Ridge aprovecha de manera eficaz las señales sectoriales y los rezagos anuales observables (\(\text{RMSE}=0{,}503\), \(R^2_{\text{OOS}}=+0{,}094\)), superando con significancia estadística a la corrección por media histórica (\(p=0{,}0071\)). Sin embargo, al incluir las variaciones extremas de 2020 y 2021 (\(-40{,}5\)\% y \(+58{,}5\)\%), los modelos lineales sin acotamiento extrapolan fuera de rango y elevan fuertemente su error en el bloque final (\(\text{RMSE}=0{,}731\)). En cambio, las técnicas clásicas acotadas evitan ese colapso: las características aleatorias de Fourier (RFF) alcanzan la mayor precisión global (\(\text{RMSE}=0{,}456\), \(R^2_{\text{OOS}}=+0{,}149\)) y los árboles potenciados acotados (GBDT) obtienen el menor error en el tramo pandémico (\(\text{RMSE}=0{,}349\), \(R^2_{\text{OOS}}=+0{,}087\)).

En tercer lugar, al evaluar el aporte específico del reservorio cuántico (QRC), los resultados muestran que su ventaja frente a la predicción nula o la corrección media clásica (\(\text{RMSE}=0{,}478\) global y \(0{,}363\) en el bloque final; \(p=0{,}0175\) sin ajuste frente a la media histórica) se explica por el acotamiento natural de sus variables de salida, mas no por superioridad cuántica intrínseca. Frente a referentes clásicos no lineales de igual dimensión de lectura, QRC empata estadísticamente con la red clásica de estado eco ESN (\(p=0{,}730\)), queda por debajo de RFF (\(p=0{,}0146\) global sin ajuste y \(p_{\text{Holm}}=0{,}0265\) antes de la pandemia) y presenta mayor variabilidad entre semillas aleatorias. Tal como preveía el diseño, este resultado negativo controlado constituye evidencia empírica valiosa para evitar atribuir ventajas injustificadas a modelos cuánticos en series macroeconómicas cortas.

En cuarto lugar, el experimento sintético calibrado y el análisis de potencia ayudan a entender por qué no surge una ventaja cuántica clara. La simulación sobre doce escenarios revela una región favorable vacía (\(\mathcal{A}=\emptyset\)) incluso al ampliar la muestra hasta \(1\,400\) períodos, debido a que la codificación angular no lineal distorsiona la señal cuando la persistencia lineal \(\phi_{12}\) es alta (\(-20{,}90\)\% frente a Ridge). Además, el análisis de potencia sobre las diferencias observadas en la muestra peruana (\(n=250\)) indica que se requerirían reducciones de error de entre \(5{,}7\)\% y \(11{,}1\)\% frente a controles clásicos acotados (\(36{,}1\)\% frente a Ridge) para lograr una potencia del 80\%, lo cual explica la amplitud del conjunto de confianza de modelos global y muestra que las series disponibles no permiten discriminar ganancias pequeñas del 2\%.

Por último, el análisis de sensibilidad confirmó la estabilidad cualitativa del sesgo positivo entre \(H=11\) y los horizontes alternativos (\(H=12, 6, 3\)) y mostró que la no linealidad detectada en la prueba RESET proviene del episodio pandémico. Cabe delimitar tres restricciones objetivas del estudio: primero, la base trabaja en régimen mensual aproximado al tiempo real ante la falta de fecha calendario diaria y el inicio posterior de los sectores de electricidad y servicios; segundo, los modelos de árboles y el circuito de cinco qubits se ejecutaron mediante simulaciones matemáticas exactas en Python sin ruido de hardware físico; y tercero, al no verificarse superioridad predictiva estable sobre los controles clásicos, se mantuvo detenida la evaluación económica condicionada sin generar beneficios monetarios hipotéticos.
